\documentclass[11pt,a4paper]{article}
\sloppy
% ==========================================
% SETTINGS.TEX - ZENTRALE KONFIGURATION
% ==========================================
\usepackage[utf8]{inputenc}
\usepackage[T1]{fontenc}
\usepackage[american]{babel}
\usepackage[top=1.8cm, bottom=1.8cm, left=1.8cm, right=1.8cm]{geometry}
\usepackage{xcolor}
\usepackage{hyperref}
\usepackage{enumitem}
\usepackage{graphicx}
\usepackage{calc}
\usepackage{charter}
\usepackage{microtype}
\usepackage{etoolbox}

% --- DESIGN FARBEN ---
\definecolor{primary}{RGB}{15, 23, 42}      % Deep Slate / Tech Dark
\definecolor{accent}{RGB}{37, 99, 235}     % MP Blue Accent
\definecolor{textgray}{RGB}{51, 65, 85}    % Charcoal Body Text

% --- PERMANENTE ABSENDER-DATEN ---
\newcommand{\meinName}{Markus Gruber, M.Sc.}
\newcommand{\meineStrasse}{Maximilianstraße 55 Stg. B 20}
\newcommand{\meinOrt}{3100 St. Pölten}
\newcommand{\meineEmail}{Markus.Gruber4@gmx.net}
\newcommand{\meineTel}{+43 650 5871145}
\newcommand{\meinLinkedin}{https://www.linkedin.com/in/markus-gruber-70851428/}
\newcommand{\meinGithub}{https://github.com/CodeByMaxx}

% --- ROLLEN-DEFAULTS ---
\providecommand{\cvRolle}{}
\renewcommand{\cvRolle}{AI & Automation Engineer -- M&A Deal Automation}
\providecommand{\cvSubheader}{}
\renewcommand{\cvSubheader}{Python \cdot LLM APIs (OpenAI, Anthropic) \cdot Financial Data Pipelines \cdot Agent Workflows \cdot RAG}

% --- EMPFÄNGER-DATEN ---
\newcommand{\firmenName}{MP Corporate Finance}
\newcommand{\ansprechPartner}{Katharina Schwarzl-Moser (HR Manager)}
\newcommand{\firmenStrasse}{Wien, Österreich}
\newcommand{\firmenOrt}{Wien, Österreich}

% --- INTELLIGENTE ADRESS-AUSGABE (Verhindert doppelte Zeilen, wenn Straße == Ort) ---
\newcommand{\firmenAdresse}{%
    \def\tempStrasse{Wien, Österreich}%
    \def\tempOrt{Wien, Österreich}%
    \ifx\tempStrasse\tempOrt
        \tempOrt%
    \else
        Wien, Österreich \\[0.15em] Wien, Österreich%
    \fi
}

\newcommand{\stellenTitel}{AI & Automation Engineer (f/m/x)}
\newcommand{\referenzNummer}{}

\newcommand{\anrede}{Dear Ms. Schwarzl-Moser,}
\newcommand{\firmenHighlight}{MP Corporate Finance is one of Europe's leading mid-market M&A boutiques, where your AI tools will be deployed on live cross-border transactions from day one.}

% --- HEADER BEFEHL ---
\newcommand{\cvHeader}{%
\begin{center}
    {\Huge\bfseries\color{primary} \meinName}\\[0.1cm]
    {\large\bfseries \cvRolle}\\[0.05cm]
    {\color{accent}\textbf{\cvSubheader}}\\[0.1cm]
    {\small \meinOrt, Austria \quad$\cdot$\quad \meineTel \quad$\cdot$\quad \href{mailto:\meineEmail}{\meineEmail} \quad$\cdot$\quad \href{\meinLinkedin}{LinkedIn} \quad$\cdot$\quad \href{\meinGithub}{GitHub}}
\end{center}
\vspace{0.1cm}
}

% --- BEFEHLE FÜR LEBENSLÄUFE & ANSCHREIBEN ---
\newcommand{\cvsection}[1]{%
    \vspace{0.25cm}%
    {\Large\bfseries\color{primary} #1}\par%
    \vspace{-0.1cm}{\color{primary}\rule{\textwidth}{0.8pt}}\par%
    \vspace{0.15cm}%
}

\newcommand{\daterange}[2]{\hfill{\small #1 -- #2}}

% --- ENGINEERING APPROACH & MINDSET DEFINITION ---
\newcommand{\engineeringApproach}{%
    \cvsection{Engineering Approach \& Mindset}%
    \begin{itemize}[leftmargin=*, itemsep=2pt]%
        \item \textbf{Systematic Problem Solving:} First-principles analytical approach to systematic debugging and stability optimization.%
        \item \textbf{Quality & Efficiency Focus:} Pragmatic engineering mindset balancing rapid deployment with long-term maintainability and high performance.%
        \item \textbf{Cross-Functional Integration:} Acts as the technical glue bridging development, QA, and architecture.%
    \end{itemize}%
}

% --- HYPERREF SETUP ---
\hypersetup{
    colorlinks=true,
    linkcolor=primary,
    urlcolor=primary,
    pdfauthor={\meinName},
    pdftitle={Application Documents - \meinName}
}


\pagestyle{empty}
\setlength{\parindent}{0pt}
\setlength{\parskip}{0.5em}

\begin{document}

% --- HEADER / SENDER ---
\begin{minipage}[t]{0.55\textwidth}
    {\Large \bfseries \color{primary} \meinName}\\[0.2em]
    {\small \color{accent} \bfseries \stellenTitel}
\end{minipage}
\begin{minipage}[t]{0.45\textwidth}
    \raggedleft \footnotesize \color{textgray}
    \meineStrasse $\cdot$ \meinOrt, Österreich\\
    E-Mail: \href{mailto:\meineEmail}{\meineEmail}\\
    Telefon: \meineTel $\cdot$ \href{\meinLinkedin}{LinkedIn} $\cdot$ \href{\meinGithub}{GitHub}
\end{minipage}

\vspace{0.8em}
\color{accent}\hrule height 1pt
\vspace{1.0em}

% --- RECIPIENT & DATE ---
\begin{minipage}[t]{0.55\textwidth}
    \color{primary} \small
    \textbf{\firmenName}\\[0.15em]
    \ansprechPartner\\[0.15em]
    \firmenAdresse\\[0.15em]
\end{minipage}
\begin{minipage}[t]{0.45\textwidth}
    \raggedleft \small \color{textgray}
    \meinOrt, \today
\end{minipage}

\vspace{1.2em}

% --- BETREFF UND DYNAMISCHER REFERENCE CODE ---
\vspace{0.8cm}
{\Large\bfseries\color{primary} Bewerbung als \stellenTitel}

\ifdefempty{\referenzNummer}{%
  \vspace{0.4cm}%
}{%
  \vspace{0.1cm}%
  {\small\color{textgray}\textbf{Referenznummer:} \referenzNummer}\\[0.4cm]%
}

\vspace{0.6em}

% --- COVER LETTER BODY ---
\color{primary} \small

\anrede

Die Verknüpfung fortgeschrittener mathematischer Modellierung mit hochperformanter, produktionsreifer KI-Infrastruktur erfordert weit mehr als bloßes Modelltraining. Sie verlangt ein tiefes numerisches Verständnis, hardwarebeschleunigte Ausführung und skalierbare Datenarchitekturen für komplexe Enterprise-Datenströme. Mit einem M.Sc. in Scientific Computing und über 9 Jahren Erfahrung in Software- und ML-Engineering bin ich darauf spezialisiert, KI-Konzepte in performante Produktionssysteme zu überführen.

Von eingebetteter GPU-Hardwarebeschleunigung (CUDA/ONNX) bis hin zu cloud-nativen RAG-Verarbeitungsbackends (Vektordatenbanken, PyTorch, AWS/Azure) bringe ich die technische Tiefe mit, um datengetriebene Intelligenz für \textbf{\firmenName} voranzutreiben.

\textbf{Warum mein Profil zu Ihrer AI- & Data-Engineering-Rolle passt:}

\begin{itemize}[leftmargin=1.2em, itemsep=0.15em, parsep=0pt, topsep=0.2em]
    \item \textbf{Mathematisches Fundament:} Master-Abschluss der Universität Wien in Scientific Computing mit fundiertem Fachwissen in numerischen Algorithmen, linearer Algebra und Datenoptimierung.
    \item \textbf{Production AI \& RAG Architectures:} Praktische Erfahrung in Konzeption und Betrieb von Cloud-basierten Retrieval-Augmented Generation (RAG) Backends, Vektordatenbanken und automatisierten PyTorch-Pipelines.
    \item \textbf{Hardware-Beschleunigung \& Datensysteme:} Nachweisbare Erfolge bei der Optimierung von Deep-Learning-Modellen mittels CUDA für Echtzeit-Edge-Processing und verteilte Daten-Pipelines (PySpark, DuckDB).
\end{itemize}

\firmenFlussbezug spricht mich insbesondere wegen \textit{\firmenHighlight} an. Gerne möchte ich meinen pragmatischen Ansatz zur Problemlösung in Ihr Engineering-Team einbringen, um komplexe Systeme nachhaltig zu verbessern.

Über die Gelegenheit, Ihre technischen Herausforderungen und Projekte in einem persönlichen Gespräch zu besprechen, freue ich mich.

\vspace{0.6em}

Mit freundlichen Grüßen

\vspace{1.2em}

\textbf{\meinName}

\end{document}
